\subsection{Realized descriptors: extrapolation and bounded alternatives}
\label{sec:bounded_results}
The historical follow-up in Section~\ref{sec:bounded_method} first reveals a
specific extrapolation failure. All historical networks are trees, so
$B9^{\rm M}=|\mathcal L|/N=1-1/N$. When the low-criticality family is excluded,
training cases contain 30--170 zones, compared with 4--30 in the omitted
family. The training mean and standard deviation of $B9^{\rm M}$ are
0.98940 and 0.00443. A four-zone case has $B9^{\rm M}=0.75$, or approximately
$-54.0$ training standard deviations. For the service preference at five
seconds, its B9 contribution is $-7.08793$, yielding an affine prediction
of $-5.94420$ against an observed target of 0.04046. Positive coefficients
do not prevent a negative prediction when centered features extrapolate.

Across 100 omitted low-criticality scenarios, five deadlines and five
preferences, 2,216 of 2,500 realized-descriptor linear predictions are
negative. B9 supplies the most negative additive contribution in 1,898 of
these cases. These are repeated predictions on 100 scenarios, not 2,500
independent instances. The explanation is an exact decomposition of the
fitted predictor, not evidence that B9 physically causes solver difficulty.

Projecting both representations' saved linear predictions onto $[0,1]$
reverses the aggregate comparison: realized descriptors reduce family-held-out
MSE by 18.6--51.3\% relative to projected generator descriptors, depending on
preference. Thus the very large unprojected losses should not be interpreted
as a general failure of the realized feature set. Projection nevertheless
collapses many easy-case predictions to zero and leaves other transfer errors.

\begin{table}[htbp]\centering\small
\begin{tabular}{lrrrr}\toprule
Preference & \shortstack{Projected\\linear} & \shortstack{Linear\\without B9} & \shortstack{Fractional\\logit} & \shortstack{Fractional logit\\without B9}\\\midrule
\multicolumn{5}{l}{\textit{Scenario-grouped validation}}\\
Equal & 0.01337 & 0.01401 & 0.00992 & 0.00992\\
Cost & 0.02313 & 0.02386 & 0.01605 & 0.01604\\
CO$_2$ & 0.01705 & 0.01748 & 0.01317 & 0.01317\\
Time & 0.00828 & 0.00899 & 0.00660 & 0.00670\\
Service & 0.01694 & 0.01757 & 0.01253 & 0.01253\\
\addlinespace
\multicolumn{5}{l}{\textit{Family-held-out validation}}\\
Equal & 0.03605 & 0.05384 & 0.01504 & 0.01451\\
Cost & 0.06047 & 0.08298 & 0.02529 & 0.02491\\
CO$_2$ & 0.03376 & 0.05674 & 0.02437 & 0.01943\\
Time & 0.01774 & 0.03702 & 0.01314 & 0.01099\\
Service & 0.04410 & 0.06225 & 0.01776 & 0.01730\\
\bottomrule
\end{tabular}
\caption{Exploratory historical MSE using realized MILP descriptors. Each cell averages 3,000 observations (600 scenarios, five deadlines). Both linear columns deliver predictions projected onto $[0,1]$; fractional-logistic outputs are bounded by their link. Models and tuning are specific to each deadline. The independent campaign contributes no outcomes.}\label{tab:bounded_milp}
\end{table}


The fractional-logistic model with B9 removed reduces MSE relative to the
projected realized-descriptor comparator by 19.1--30.6\% under scenario folds
and by 38.1--60.8\% under family holdout (Table~\ref{tab:bounded_milp}). All
five pooled contrasts in each scheme have positive lower bounds in the
conditional simultaneous intervals. Removing B9 from the linear model alone
instead increases MSE by 2.5--8.6\% and 37.2--108.7\%, respectively. The
choice of predictive model therefore changes the interpretation of ablation:
B9 can be useful to the linear approximation even though its extrapolation
requires control.

Within the bounded model, B9 removal further improves all five family-held-out
losses, with positive conditional simultaneous lower bounds. The additional
MSE reductions range from 1.5\% to 20.3\% relative to the full bounded model.
Under scenario folds, effects are negligible for four preferences, and the
time-priority MSE increases by 1.5\%. The historical transfer evidence favors
the bounded model without B9 as a candidate, without establishing a universal
advantage for dropping that coordinate.

The original generator descriptors also benefit from fractional-logistic
fitting (Appendix~\ref{app:bounded_followup}). With the same bounded model
and B9 ablation, realized descriptors have lower MSE for all five preferences
under scenario folds and for four under family holdout. The exception is
CO$_2$ priority, whose family-held-out MSE is 2.7\% higher. Accordingly, the
large improvements over projected linear predictions must not be attributed
entirely to descriptor recomputation.

Transfer remains imperfect. For the bounded realized model without B9,
family-held-out MSE is 1.38--1.64 times the corresponding scenario-fold MSE.
Within the low-criticality family itself, four preferences improve relative
to the projected comparator, but time-priority MSE rises from 0.01006 to
0.01293. These historical results neither establish deployment calibration
nor revise the ES-allocation results in Section~\ref{sec:revision_allocation}.
No outcome from the independent campaign is used in these fits or comparisons.
